\section{Experimental Setup}
\label{sec:experimental_setup}

\subsection{Workstation Hardware Configuration}
All empirical benchmarks are conducted on a commodity developer workstation running Ubuntu Linux 24.04 LTS, equipped with an Intel Core Ultra 5 processor, 16~GB DDR5 RAM, and a 1~TB NVMe SSD.

\subsection{Benchmark Datasets}
Workloads are synthesized using the industry-standard TPC-H benchmark suite \cite{tpch2021} generated via embedded DuckDB engines across three distinct Scale Factors:
\begin{itemize}
    \item \textbf{SF-1} ($1$~GB raw relational tables)
    \item \textbf{SF-5} ($5$~GB raw relational tables)
    \item \textbf{SF-10} ($10$~GB raw relational tables)
\end{itemize}
All tables are stored as compressed Parquet files with Zstandard (ZSTD) compression to reflect modern cloud-native edge storage practices.

\subsection{Evaluation Metrics}
System performance is evaluated along three primary dimensions:
\begin{enumerate}
    \item \textbf{Peak Resident Set Size (RSS):} Continuous high-frequency memory tracking sampled at $50$~ms intervals via OS \texttt{psutil} daemons.
    \item \textbf{Token Consumption:} Cumulative prompt and completion tokens exchanged with the remote reasoning gateway.
    \item \textbf{Task Execution Success Rate:} Successful completion without unhandled runtime exceptions or out-of-memory kernel panics.
\end{enumerate}
